\documentclass[11pt,aspectratio=169]{beamer}
%%%%%%%%% GENERAL PACKAGES
%\usepackage[dvipsnames]{xcolor}
%\usepackage{pdfpages}
%\usetheme[progressbar=frametitle]{metropolis}
%\setbeamercolor{background canvas}{bg=white}
%\usepackage{appendixnumberbeamer}
%\usepackage{booktabs}
%\usepackage[scale=2]{ccicons}
%\usepackage{pgfplots}
%\usepgfplotslibrary{dateplot}
%\usepackage{xspace}
%\newcommand{\themename}{\textbf{\textsc{metropolis}}\xspace}
%\usepackage[absolute,overlay]{textpos}






%%%%%%%%% COLOR THEME

% Define some colors:
\definecolor{DarkFern}{HTML}{407428}
\definecolor{DarkCharcoal}{HTML}{4D4944}
\definecolor{AlertColor}{RGB}{89,124,158}
\definecolor{HighLight}{RGB}{96,95,134}
\definecolor{Important}{RGB}{234,122,133}
\definecolor{Yellow}{HTML}{00539C}
\colorlet{Fern}{DarkFern!85!white}
\colorlet{Charcoal}{DarkCharcoal!85!white}
\colorlet{LightCharcoal}{Charcoal!50!white}
\colorlet{HighLight2}{AlertColor}
\colorlet{DarkRed}{red!70!black}
\colorlet{DarkBlue}{blue!70!black}
\colorlet{DarkGreen}{green!70!black}
\definecolor{RoyalBlue}{HTML}{00539C}
\definecolor{Peach}{HTML}{EEA47F}
\definecolor{ForestGreen}{HTML}{2C5F2D}
\definecolor{MossGreen}{HTML}{E8FCC9}
\definecolor{SeaGreen}{HTML}{2E8B57}
% Use the colors:
\setbeamercolor{title}{fg=Fern}
\setbeamercolor{frametitle}{fg=MossGreen,bg=ForestGreen}
\setbeamercolor{normal text}{fg=Charcoal!70!black}
\setbeamercolor{block title}{fg=black,bg=Fern!25!white}
\setbeamercolor{block body}{fg=black,bg=Fern!10!white}
\setbeamercolor{block title alerted}{fg=black,bg=DarkRed!25!white}
\setbeamercolor{block body alerted}{fg=black,bg=DarkRed!10!white}
\setbeamercolor{alerted text}{fg=DarkRed}
\setbeamercolor{itemize item}{fg=Charcoal}



%%%%%%%%% OTHER COMMANDS
\newcommand{\indep}{\perp\!\!\! \perp}
\newcommand{\comment}[1]{}
\newcommand{\bs}{\boldsymbol}
\newcommand{\tr}{\text{trace}}
\newcommand{\sgn}{{\rm sgn}}
\def\T{\top}
%\newcommand{\det}{\text{det}}
\newcommand{\var}{\mathrm{var}}
\newcommand{\cC}{{\cal C}}
\renewcommand{\d}{{\rm d}}
\newcommand{\cG}{{\cal G}}
\newcommand{\cV}{{\cal V}}
\newcommand{\cE}{{\cal E}}
\newcommand{\cM}{{\cal M}}
\newcommand{\cP}{{\cal P}}
\newcommand{\cX}{{\cal X}}
\newcommand{\cY}{{\cal Y}}
\newcommand{\X}{\mathbf{X}}
\newcommand{\Y}{\mathbf{Y}}
\newcommand{\x}{\mathbf{x}}
\newcommand{\y}{\mathbf{y}}
\newcommand{\z}{\mathbf{z}}

\newcommand{\argmin}{\operatornamewithlimits{argmin}}
\newcommand{\eps}{\varepsilon}
\newcommand{\<}{\langle}
\renewcommand{\>}{\rangle}


%

\setbeamertemplate{navigation symbols}{}
\setbeamertemplate{footline}[text line]{%
    \hfill\strut{%
        \scriptsize\sf\color{black!60}%
        \quad\insertframenumber/\inserttotalframenumber
    }
    %\hfill
    }


\usenavigationsymbolstemplate{}
\setbeamersize{text margin left=.2cm,text margin right=.2cm} 
\addtobeamertemplate{frametitle}{}{\vspace{-1.2mm}}
\setbeamertemplate{itemize item}{$\bullet$}

\setbeamertemplate{itemize subitem}{\tiny\raise1.5pt\hbox{\donotcoloroutermaths$\blacktriangleright$}}
\setbeamertemplate{itemize subsubitem}{\tiny\raise1.5pt\hbox{\donotcoloroutermaths$\blacktriangleright$}}
\setbeamertemplate{enumerate item}{\insertenumlabel.}
\setbeamertemplate{enumerate subitem}{\insertenumlabel.\insertsubenumlabel}
\setbeamertemplate{enumerate subsubitem}{\insertenumlabel.\insertsubenumlabel.\insertsubsubenumlabel}
\setbeamertemplate{enumerate mini template}{\insertenumlabel}






\newcommand{\TODO}[1]{{\color{red}{[TODO: #1]}}}


\newcommand{\R}{\mathbb R}
\newcommand{\E}{\mathbb E}
\renewcommand{\P}{\mathbb P}


\DeclareMathOperator*{\cov}{cov}


\newsavebox{\zerobox}
\newenvironment{nospace}
{\par\edef\theprevdepth{\the\prevdepth}\nointerlineskip
  \setbox\zerobox=\vtop to 0pt\bgroup
  \hrule height0pt\kern\dimexpr\baselineskip-\topskip\relax
}
{\par\vss\egroup\ht\zerobox=0pt \wd\zerobox=0pt \dp\zerobox=0pt
  \box\zerobox}

\usepackage{soul}
\makeatletter
\let\HL\hl
\renewcommand\hl{%
  \let\set@color\beamerorig@set@color
  \let\reset@color\beamerorig@reset@color
  \HL}
  \makeatother


\title[Calculus and Linear Algebra]{Lecture 5: Calculus and Linear Algebra}
\author[Piotr Zwiernik, Barcelona School of Economics]{Piotr Zwiernik\\[0.5ex]Mathematics Brush-up}
\titlegraphic{\includegraphics[width=1.5in]{img/bse.png}}
\date{}

\begin{document}

% ---------------- Title ----------------
\begin{frame}
  \titlepage
\end{frame}

% ---------------- Slide 1 ----------------
\begin{frame}{Motivation: why separable ODEs?}
\begin{small}
Ordinary differential equations (ODEs) appear whenever we model how a quantity changes over time or another variable.  

\medskip
A simple but very important class: \textcolor{SeaGreen}{separable ODEs}, where variables can be separated:
\[
\frac{dy}{dx} = g(x)\,h(y).
\]

\begin{itemize}
  \item Widely used in economics, biology, and physics for growth/decay models.
  \item Method: rearrange into \(\frac{dy}{h(y)} = g(x)\,dx\), then integrate both sides.
  \item Today: a very small illustrative example.
\end{itemize}
\end{small}
\end{frame}

% ---------------- Slide 2 ----------------
\begin{frame}{A toy separable ODE}
\begin{small}
Consider
\[
v'(x)\;=\;\frac{{\rm d}v}{{\rm d}x}\;=\;-\frac{1}{x}\,v(x).
\]

\medskip
\textcolor{SeaGreen}{Step 1: Separate variables.}
\[
\frac{dv}{v} = -\frac{dx}{x}.
\]

\textcolor{SeaGreen}{Step 2: Integrate both sides.}
\[
\int \frac{dv}{v} = -\int \frac{dx}{x}
\quad\Rightarrow\quad
\ln|v(x)|=-\ln|x|+C.
\]

Hence
\[
v(x)=C_1\,x^{-1}, \quad x\neq 0.
\]
\end{small}
\end{frame}

% ---------------- Slide 3 ----------------
\begin{frame}{Recovering $u$ from $v=u'$}
\begin{small}
Suppose \(v(x)=u'(x)\). Then
\[
u'(x)=C_1\,x^{-1}.
\]

\textcolor{SeaGreen}{Integrate again:}
\[
u(x)=C_1\ln|x|+C_2.
\]

\medskip
\textbf{Domain note:} Because of $\ln|x|$, we must restrict to $x>0$ or $x<0$ separately.
\end{small}
\end{frame}

% ---------------- Slide 4 ----------------
\begin{frame}{Example with initial conditions}
\begin{small}
Suppose $u(1)=0$ and $u'(1)=2$.

\begin{align*}
u'(1)=2 &\;\Rightarrow\; C_1=2, \\
u(1)=0 &\;\Rightarrow\; 0=2\ln 1+C_2 \;\Rightarrow\; C_2=0.
\end{align*}

Therefore
\[
u(x)=2\ln x \quad \text{(valid on $x>0$)}.
\]
\end{small}
\end{frame}

% ---------------- Slide 5 ----------------
\begin{frame}{Key takeaways}
\begin{small}
\begin{itemize}
  \item A \textcolor{SeaGreen}{separable ODE} can be solved by separating $y$-terms and $x$-terms, then integrating.
  \item Constants of integration appear at each step and are fixed by initial conditions.
  \item Always check the domain: functions like $\ln|x|$ force $x>0$ or $x<0$.
  \item Even this tiny example illustrates the general workflow.
\end{itemize}
\end{small}
\end{frame}
\end{document}